\documentclass[tikz,border=6pt]{standalone}
\usepackage{peabra}
\begin{document}
\begin{tikzpicture}[x=1cm,y=1cm, s/.style={doboz=50mm, font=\footnotesize, minimum height=10mm}]
  \node[sotetdoboz=50mm, font=\footnotesize, minimum height=10mm] (d0) at (0,0) {a döntési igény megértése a rendszer összefüggésében};
  \node[s] (d1) at (0,-1.5) {értékelési kritériumok meghatározása (kötelező és javító)};
  \node[s] (d2) at (0,-3.0) {alternatívák azonosítása (döntési fa a szűkítéshez)};
  \node[s] (d3) at (0,-4.5) {értékelési módszer és eszköz kiválasztása};
  \node[s] (d4) at (0,-6.0) {az alternatívák értékelése a kritériumok szerint};
  \node[zolddoboz=50mm, font=\footnotesize, minimum height=10mm] (u) at (0,-7.5) {megváltoztathatja-e a rangsort a bizonytalanság csökkentése?};
  \node[kiemeltdoboz=50mm, font=\footnotesize, minimum height=10mm] (d5) at (0,-9.0) {javaslat a rangsor robusztusságával együtt};
  \node[s] (d6) at (0,-10.5) {eredmények jelentése, munkatermékek rögzítése};
  \foreach \a/\b in {d0/d1,d1/d2,d2/d3,d3/d4,d4/u,d5/d6} \draw[nyil] (\a) -- (\b);
  \draw[nyil] (u) -- node[felirat, right] {nem} (d5);
  \draw[nyil] (u.east) -- ++(1.0,0) |- node[felirat, pos=0.25, right, align=left] {igen, és megéri:\\több adat,\\újabb elemzés} (d4.east);
  \draw[nyil, draw=peSzurke, dash pattern=on 3pt off 2pt] (d5.west) -- ++(-1.0,0) |- node[felirat, pos=0.25, left, align=right] {ha alacsonyabb\\pontszámú\\alternatívát\\ajánlanánk:\\a kritériumok\\újragondolása} (d1.west);
\end{tikzpicture}
\end{document}
